% =============================================================
%  Chapter — Harmonic and high-frequency interference
%  This file holds the chapter heading and its introduction only; the
%  content is one file per subsection, \input below in reading order.
%    regulatory-basis.tex       US regulatory status of PLC and RI, how the
%                               limits are set, adaptations to the western US
%    plc.tex                    power line carrier band: coupling paths,
%                               form of the requirement, indicative values
%    telephone-interference.tex equivalent disturbing current of the DC line
%    radio-interference.tex     RI of the converter station: sources, limits,
%                               measurement conventions, altitude correction
%    dc-line-corona.tex         corona RI and audible noise of the DC line
%    verification.tex           design, factory and site verification
%  Sources: [27] 47 CFR Part 15; [28] UTC PLC database; [29] CIGRE TB 139;
%           [30] CIGRE TB 553; [31] CISPR TR 18-2; [32] CISPR 16-1-1;
%           [33] ITU-T O.41; [34] IEEE 1124; [35] CISPR 11; [36] NEMA 107;
%           [37] EPRI HVDC Transmission Line Reference Book; [38] IEEE 430.
%  The chapter has no figures yet; fig/ is reserved for them.
% =============================================================
\section{Harmonic and High-Frequency Interference}
\label{sec:interference}

A VSC HVDC converter station emits conducted and radiated disturbances from
the fundamental frequency up to the VHF range. The low-order harmonics below
3\,kHz are the subject of the harmonic performance specification of a
project and are not treated here. This chapter covers the bands above them:
the power line carrier (PLC) band from 30\,kHz to 500\,kHz, where converter
noise disturbs the carrier channels that utilities run over their
transmission lines, and the radio interference (RI) band from 150\,kHz to
30\,MHz, where corona and converter switching disturb broadcast and radio
services. Table~\ref{tab:interference-bands} lists the bands with their
dominant source and coupling path. For each band the chapter gives the
regulatory position in the United States, the form in which a project
specification states the requirement, the measurement conventions that make
the requirement verifiable, and indicative values with their basis.

\begin{table}[!ht]
  \centering
  \renewcommand{\arraystretch}{1.25}
  \begin{tabularx}{\linewidth}{@{}l l >{\raggedright\arraybackslash}X >{\raggedright\arraybackslash}X l@{}}
    \toprule
    \textbf{Band} & \multicolumn{1}{c}{\textbf{Frequency}} & \textbf{Dominant source} & \textbf{Coupling path} & \multicolumn{1}{c@{}}{\textbf{Section}} \\
    \midrule
    Harmonics & 60\,Hz -- 3\,kHz & Converter switching & Conducted, AC and DC terminals & --- \\
    Telephone interference & 60\,Hz -- 5\,kHz & Harmonic currents on the DC line & Inductive coupling to telecommunication circuits & \ref{sec:tif} \\
    PLC band & 30\,kHz -- 500\,kHz & Converter switching & Conducted on the AC and DC lines; inductive coupling to parallel lines & \ref{sec:plc} \\
    Radio interference & 150\,kHz -- 30\,MHz & Converter switching; corona on high-voltage apparatus & Radiated from valve hall, buswork and switchyards & \ref{sec:ri} \\
    VHF/UHF interference & 30\,MHz -- 1\,GHz & Converter switching; gap discharges & Radiated & \ref{sec:ri} \\
    \bottomrule
  \end{tabularx}
  \caption{Emission bands of a converter station, their dominant source and coupling path, and the section that treats them.}
  \label{tab:interference-bands}
\end{table}

\subsection{Regulatory Basis and the Form of the Requirement}
\label{sec:interference-basis}

In the United States no numerical emission limit applies to high-voltage
transmission equipment in the PLC and RI bands. Overhead lines and converter
stations are incidental radiators under 47~CFR Part~15, \S\,15.5 and
\S\,15.13~[27]: the operator is obliged not to cause harmful interference to
licensed radio services and to apply good engineering practice, and must
correct interference once it occurs, but no field strength or noise voltage
is prescribed. PLC systems themselves operate unlicensed under \S\,15.113 in
the band 9\,kHz to 490\,kHz on the utility's own lines; their frequency
assignments are coordinated between utilities through the PLC database of
the Utilities Technology Council (UTC)~[28]. In Europe the position is the
same in substance, with CISPR TR~18~[31] as the technical basis and no binding
limit for high-voltage installations.

The limits in a project specification are therefore design values set by the
owner. The methodology for converter-generated noise is that of CIGRE
TB~139~[29] and its successor TB~553~[30], which treat PLC and RI noise as
part of the AC filter design; the methodology for radio interference from
high-voltage installations, including the measurement procedure and the
derivation of limits from the protection of broadcast reception, is that of
CISPR TR~18-2~[31]. A clause written on this basis has five elements: the
frequency range, the limit, the measuring point and distance, the weather
condition and statistic, and the detector and bandwidth. A clause that omits
one of them cannot be verified against the supplier's design or on site.

Three conditions of the western United States change the values relative to
a European or eastern-US template. The corona onset of high-voltage hardware
falls with air density: at 1500\,m the relative air density is 0.84, and a
bundle or a fitting that meets a corona criterion at sea level does not meet
it at that elevation, so every corona-related value is calculated and
verified at the site or route elevation. The climate is dry, and for DC
corona the fair-weather emission of the positive pole is the governing case,
so the limits are stated for fair weather rather than for the foul-weather
case that AC line specifications emphasise. The population along a route is
sparse, with a small number of communities, while the interconnection areas
carry a dense set of third-party transmission lines on which carrier
protection remains in service; the PLC requirement is therefore written to
protect those third-party carriers, and the RI requirement is referred to the
station boundary and the edge of the right-of-way rather than to the nearest
dwelling.
\FloatBarrier

\subsection{Power Line Carrier Band}
\label{sec:plc}

\subsubsection{Sources and coupling paths}
Converter noise in the PLC band originates in the switching of the valves. In
a line-commutated converter the commutation notches excite the stray
capacitances and inductances of the valve hall and buswork into oscillations
in the range of tens to hundreds of kilohertz; in a modular multilevel
converter the individual submodule switchings are small steps, and the
emission in this band is lower by one to two orders of magnitude, but it is
not zero and still depends on the switching frequency of the control and on
the layout of the valve hall. The noise leaves the station on three paths: as
a conducted current on the AC lines through the converter transformers and
the AC switchyard, as a conducted current on the DC line, and by inductive
coupling from the DC line into any transmission line that runs parallel to or
crosses it. The third path is the one a station-terminal measurement does not
capture.

\subsubsection{Form of the requirement}
A PLC clause has two parts. The first is an absolute limit on the noise
injected into each connected line, stated as a noise voltage in a 3\,kHz
bandwidth, the channel bandwidth of a carrier receiver, measured with a
quasi-peak detector according to CISPR~16-1-1~[32] at the line side of the
switchyard. The second is a signal-to-noise criterion for each carrier
frequency in service on a protected line: the noise coupled into that line
may not reduce the carrier signal-to-noise ratio at the receiver below a
stated margin, with the minimum received carrier level and the receiver
bandwidth given by the line owner and the coupling calculated from the
geometry of the parallel or crossing section. The second part governs
wherever it is stricter. The list of protected lines and carrier frequencies
is compiled by the owner from the UTC database~[28] and from the owners of the
lines within a stated corridor of the route, and is handed to the supplier
with the DC line data. The supplier's response is the PLC filter on the AC
and DC side, tuned to the protected frequencies, and a demonstration that the
limits hold with any redundant filter branch out of service and, in a dual
symmetrical monopole, with one converter out of service.

\subsubsection{Indicative values}
A margin of 20\,dB on the carrier signal-to-noise ratio is the value in common
use for teleprotection channels. For the absolute limit, 60\,dB($\mu$V),
equal to 1\,mV in 3\,kHz, is representative of the values specified for
line-commutated schemes~[29]; whether it protects a given carrier channel
follows from the signal-to-noise criterion, which is why that criterion
governs. A VSC station meets it with a wide margin; the clause is nevertheless kept so
that the supplier must demonstrate it. A corridor of 2\,km for the survey of
protected lines covers the distance within which inductive coupling from a
DC line is calculable above the natural noise of the coupled line. For a
private network that is not connected to the wider AC grid and runs its
teleprotection over fibre, the PLC clause protects third-party lines through
the DC route only, and the AC-side part reserves carrier capability for a
later buildout.
\FloatBarrier

\subsection{Telephone Interference from the DC Line}
\label{sec:tif}

Harmonic currents on the DC line induce voltages in telecommunication circuits
that share the route. The measure is the equivalent disturbing current
$I_{eq}$: the harmonic currents of both poles, weighted with the psophometric
curve of ITU-T O.41~[33], which represents the sensitivity of the ear and the
telephone receiver, and referred to 800\,Hz, where the weighting is unity.
IEEE Std~1124~[34] gives the calculation along the line for the different
converter configurations. A limit of 500\,mA at any point of the line is the
value in common use; it was set for line-commutated schemes with open-wire
telephone circuits, and a VSC scheme reaches a small fraction of it. The
configurations to be covered are all those a scheme can operate in, since the
earth-return and metallic-return currents of a bipole in monopolar operation
and the unbalance of a dual symmetrical monopole with one converter out of
service give the largest values.
\FloatBarrier

\subsection{Radio Interference of the Converter Station}
\label{sec:ri}

\subsubsection{Sources}
Two sources radiate from a converter station in the band 150\,kHz to 30\,MHz.
Corona on the high-voltage apparatus of the AC and DC switchyards, on
bushings, connectors, corona rings and busbar hardware, has the spectrum of
an AC or DC line and falls with frequency, so that it is characterised at a
reference frequency of 0.5\,MHz (1\,MHz in older North American practice).
Valve switching radiates from the valve hall through its openings, the
wall bushings and the connections to the switchyards, with a spectrum set by
the switching edges and the resonances of the hall and buswork; it extends
beyond 30\,MHz and is the source of the VHF/UHF emission in
Table~\ref{tab:interference-bands}. Corona is controlled at the source, by
the design of the hardware for the site air density; switching emission is
controlled by the shielding of the valve hall.

\subsubsection{Form of the requirement and indicative values}
The requirement is a field-strength limit at a stated distance outside the
station fence, with a quasi-peak detector and the bandwidths of
CISPR~16-1-1~[32], in fair weather, with the converter at the operating point
of maximum emission. Table~\ref{tab:ri-limits} gives indicative values at
30\,m. In the band 150\,kHz to 30\,MHz the limit is referred to 0.5\,MHz and
follows the frequency characteristic of CISPR TR~18-2~[31] elsewhere in the
band; 50\,dB($\mu$V/m) leaves a signal-to-noise ratio
of 24\,dB against a broadcast field strength of 5\,mV/m, equal to
74\,dB($\mu$V/m), the city-grade service contour of an AM station in FCC
practice, which is within the range CISPR TR~18-2 associates with acceptable
to good reception. Above 30\,MHz the limits of
CISPR~11~[35], Group~1, Class~A, apply, since no transmission-line standard
covers that band.

\begin{table}[!ht]
  \centering
  \renewcommand{\arraystretch}{1.25}
  \begin{tabular}{@{}l c r@{}}
    \toprule
    \textbf{Frequency range} & \multicolumn{1}{c}{\textbf{Reference}} & \multicolumn{1}{c@{}}{\textbf{Limit}} \\
    \midrule
    150\,kHz -- 30\,MHz   & 0.5\,MHz, CISPR TR~18-2 characteristic & 50\,dB($\mu$V/m) \\
    30\,MHz -- 230\,MHz   & CISPR~11 Group~1 Class~A                & 30\,dB($\mu$V/m) \\
    230\,MHz -- 1\,GHz    & CISPR~11 Group~1 Class~A                & 37\,dB($\mu$V/m) \\
    \bottomrule
  \end{tabular}
  \caption{Indicative radio interference field-strength limits at 30\,m outside the station fence, quasi-peak detector, fair weather.}
  \label{tab:ri-limits}
\end{table}

\subsubsection{Apparatus and altitude}
The corona contribution is specified per item of apparatus as a radio
interference voltage (RIV) measured according to NEMA~107~[36] at a test
voltage of $1.1 \cdot U_m / \sqrt{3}$ for AC apparatus and $1.1 \cdot
U_{DC\_max}$ for DC apparatus, together with the absence of visible corona at
the same voltage. The product standards state the RIV limits at standard
atmosphere; for a station at 1500\,m the test voltage is corrected to the
relative air density of 0.84, which raises the required corona onset of the
hardware by the same factor as the creepage and clearance corrections of
Section~\ref{sec:creepage}.
\FloatBarrier

\subsection{Corona Performance of the DC Line}
\label{sec:dc-line-corona}

The DC line is not part of the converter supply, but its corona criteria are
set on the same basis as those of the DC switchyard and the station emission
at the DC terminal adds to the noise of the line, so they are given here. On a
DC line the positive pole produces the corona radio noise, and its emission
is highest in fair weather: rain lowers the DC corona noise, the opposite of
the AC case, because the space charge around the positive conductor is
dispersed by the water drops. The design statistic is therefore the fair-weather level not exceeded 50\,\% of the time, at 0.5\,MHz with a quasi-peak
detector, at a lateral distance from the positive pole, typically 15\,m from
the conductor or the edge of the right-of-way. Values of 40\,dB($\mu$V/m) to
50\,dB($\mu$V/m) at the edge of the right-of-way are the range of the design
criteria used for $\pm$500\,kV lines in the western United States~[37], chosen
to protect AM broadcast reception at the signal strengths found along the
routes. The matching audible-noise criterion is an A-weighted level of
42\,dB(A) in fair weather, 50\,\% statistic, at the same point. Both are
calculated for the maximum elevation of the route; a route through Utah and
Nevada runs between 1200\,m and 2000\,m, where the relative air density lies
between 0.87 and 0.79 and the corona onset gradient falls by the same ratio.
\FloatBarrier

\subsection{Verification}
\label{sec:interference-verification}

Compliance is demonstrated in three stages, and a specification names all
three.

At the design stage the supplier submits an interference study according to
CIGRE TB~553~[30]: the calculated noise spectrum at the AC and DC line
terminals in the band 30\,kHz to 500\,kHz, the calculated radio interference
field strength at the station fence in the band 150\,kHz to 1\,GHz, the
calculated $I_{eq}$ of the DC line, and the design of the PLC filters and the
valve-hall shielding. A preliminary study belongs to the proposal, since the
PLC filter is part of the price; the final study belongs to the design
review.

At the factory stage the RIV of the switchyard apparatus is verified by the
type tests of the product standards, with the test voltage corrected to the
site elevation as in Section~\ref{sec:ri}.

At site the background level is measured at the defined points before first
energisation and the measurement is repeated with the station energised at
0\,\%, 25\,\%, 50\,\%, 75\,\% and 100\,\% of rated power in both power
directions, in fair weather, according to IEEE Std~430~[38] for the field
strength and with coupling capacitors at the line terminals for the PLC-band
noise. The background measurement is what makes the result attributable: in
an interconnection area the natural noise of the surrounding lines can exceed
a limit on its own, and the criterion then becomes that the station raises
the measured level by not more than 3\,dB. The measurement report is part of
the site acceptance documentation.
\FloatBarrier

